% subfigure.tex - Illustrating the use of the subfigure package

% Andrew Roberts - September 2003

\documentclass[english]{article}

\usepackage{times}
\usepackage{babel}
\usepackage{graphicx}
\usepackage{subfigure}

\begin{document}

\title{Latex Tutorial 6 (Floats, Figures and Captions) Examples}

\author{Andrew Roberts}
\maketitle

\begin{figure}[htp]
  \begin{center}
    \subfigure[Original image]{\label{edge-a}\includegraphics[scale=0.75]{toucan.eps}}
    \subfigure[After Laplace edge detection]{\label{edge-b}\includegraphics[scale=0.75]{laplace_toucan.eps}}
    \subfigure[After Sobel edge detection]{\label{edge-c}\includegraphics[scale=0.75]{sobel_toucan.eps}}
  \end{center}
  \caption{Various edge detection algorithms}
  \label{edge}
\end{figure}

The Sobel algorithm gives the best results (see \ref{edge-c}).

\begin{figure}[htp]
  \begin{center}
    \subfigure[Original image]{\label{fig:edge-a}\includegraphics[scale=0.75]{toucan.eps}}
    \subfigure[After Laplace edge detection]{\label{fig:edge-b}\includegraphics[scale=0.75]{laplace_toucan.eps}} \\
    \subfigure[After Sobel edge detection]{\label{fig:edge-c}\includegraphics[scale=0.75]{sobel_toucan.eps}}
  \end{center}
  \caption{Various edge detection algorithms}
  \label{fig:edge}
\end{figure}

\end{document}
